\documentclass{beamer}

\usepackage{./packages/packages_mathalea}
\input{./competences/competences_latex.tex}
\usepackage{fontawesome}
\title{Exercices de début de cours (edc)}
\subtitle{Automatismes et autres compétences -- 1$^{\text{re}}$ année du collège}
\author{ns}
\setbeamerfont{frametitle}{size=\normalsize}
% define a new counter
\newcounter{exercice}
\setcounter{exercice}{0}
\begin{document}

%%%%%%%%%%%%%%%%%%%%%%%%%%%%%%%%%%%%%%%%%%%%%%%%%%%%%%%%%%%%%%%%%%%%%%%%%%%%%%%
%%%%%%%%%%%%%%%%%%%%%%%%%%%%%%%%%%%%%%%%%%%%%%%%%%%%%%%%%%%%%%%%%%%%%%%%%%%%%%%
%%%%%%%%%%%%%%%%%%%%%%%%%%%%%%%%%%%%%%%%%%%%%%%%%%%%%%%%%%%%%%%%%%%%%%%%%%%%%%%
\begin{frame}
\titlepage
\end{frame}

\stepcounter{exercice}
\begin{frame}{\GetCompetence{FactTri}, edc \theexercice}
	% increment the counter
 \onslide<1->{ Utiliser la méthode de résolution du deuxième degré pour factoriser l'expression suivante. 
Dans la forme factorisée, tous les coefficients doivent être entiers.
 \[15x^2+32x-7\]}
 \onslide<2->{$\Delta=32^2-4\cdot 15\cdot (-7)=1444$. Les racines du polynôme sont 
	 \[S=\left\{\dfrac{1}{5};-\dfrac{7}{3}\right\}.\]

	 La forme factorisée est 
 \[(5x-1)(3x+7)\]}
\end{frame}

\stepcounter{exercice}
\begin{frame}{ \GetCompetence{CombLin}, edc \theexercice}
	\onslide<1->{
Résoudre les systèmes d'équations suivants par combinaison linéaire :
\begin{multicols}{2}
\begin{enumerate}
	\item  $\begin{cases}\begin{aligned}-x-2y &=1\\-6x-4y &=-26\end{aligned}\end{cases}$
	\item  $\begin{cases}\begin{aligned}-4x+6y &=-44\\-3x+4y &=-32\end{aligned}\end{cases}$
\end{enumerate}
\end{multicols}}
\onslide<2->{Les solutions sont :
	\begin{multicols}{2}
\begin{enumerate}
\item $S=\left\{\left(7;-4\right)\right\}$.
\item $S=\left\{\left(8;-2\right)\right\}$.
\end{enumerate}
\end{multicols}}
\end{frame}

\stepcounter{exercice}
\begin{frame}{\GetCompetence{3f3}, edc \theexercice}
	\onslide<1->{Résoudre le système d'équations suivant dans $\mathbb{R}^3$ :
	\[\begin{cases}\begin{aligned}x+y-2z &=3\\3x+2y-z &=12\\8x-3y-6z &=-18\end{aligned}\end{cases}\]}
	\onslide<2->{La solution est $S=\left\{\left(1;\dfrac{16}{3};\dfrac{5}{3}\right)\right\}$.}	
\end{frame}

\stepcounter{exercice}
\begin{frame}{\GetCompetence{3f3}, edc \theexercice}
	\onslide<1->{Résoudre le système d'équations suivant dans $\mathbb{R}^3$ :
	\[\begin{cases}\begin{aligned}x+2y-z &=2\\x-y+2z &=5\\2x+2y+2z &=12\end{aligned}\end{cases}\]}
	\onslide<2->{La solution est $S=\left\{\left(1;2;3\right)\right\}$.}	
\end{frame}

\stepcounter{exercice}
\begin{frame}{\GetCompetence{ProbSys}, edc \theexercice}
	\onslide<1->{La recette d'une théâtre s'élève à CHF 13450. Les places sont vendues à CHF 30 et CHF 40. Sachant qu'il y a 400 places vendues, déterminer le nombre de chaque type de place.}
		\begin{itemize}
		\item<2-> $x=$ nombre de places à CHF 30 et $y=$ nombre de places à CHF 40 
		\item<3-> \[\begin{cases}\begin{aligned}x+y &=400\\30x+40y &=13450\end{aligned}\end{cases}\]
		\item<4-> $S=\left\{\left(255\,;\,145\right)\right\}$.
\end{itemize}
\end{frame}

\stepcounter{exercice}
\begin{frame}{\GetCompetence{ProbSys}, edc \theexercice}
	\onslide<1->{Un entrepreneur doit déplacer 460 tonnes de terre : il dispose de 2 camions, l'un pouvant transporter 5 tonnes et l'autre 3 tonnes ; il désire effectuer 100 transports. Combien de fois doit-il utiliser chaque camion~?}
		\begin{itemize}
		\item<2-> $x=$ nombre de transports avec le camion à 5 tonnes
		\item<2-> $y=$ nombre de transports avec le camion à 3 tonnes
		\item<3-> \[\begin{cases}\begin{aligned}x+y &=100\\ 5x+3y &=4600\end{aligned}\end{cases}\]
		\item<4-> $S=\left\{\left(80\,;\,20\right)\right\}$.
		\item<5-> Il doit utiliser 80 fois le camion à 5 tonnes et 20 fois le camion à 3 tonnes.
\end{itemize}
\end{frame}

\stepcounter{exercice}
\begin{frame}{\GetCompetence{ProbSys}, edc \theexercice}
	\onslide<1->{On doit répartir des élèves dans des groupes pour une excursion. Si on met 188 élèves par groupe, alors on au besoin de 2 groupes de moins que si on met 141 élèves par groupe. Combien d'élèves y a-t-il~?}
		\begin{itemize}
		\item<2-> $x=$ nombre total d'élèves		\item<2-> $y=$ nombre de groupes avec $141$ élèves		\item<3-> \[\begin{cases}\begin{aligned}y\cdot 141&=x\\ +(y-2)\cdot 188  &=x\end{aligned}\end{cases}\]
		\item<4-> $S=\left\{\left(1128\,;\,6\right)\right\}$.
		\item<5-> Il y a donc 1128 élèves au total.
\end{itemize}
\end{frame}

\stepcounter{exercice}
\begin{frame}{\GetCompetence{3f3}, edc \theexercice}
	\onslide<1->{Résoudre le système d'équations suivant dans $\mathbb{R}^3$ :
	\[\begin{cases}\begin{aligned}10x+3y+2z &=8\\4x+2y-5z &=-8\\3x-7y+2z &=6\end{aligned}\end{cases}\]}
	\onslide<2->{La solution est $S=\left\{\left(\dfrac{202}{447};-\dfrac{52}{447};\dfrac{856}{447}\right)\right\}$.}	
\end{frame}
\stepcounter{exercice}
\begin{frame}{\GetCompetence{Fact4}, edc \theexercice}
	\onslide<1->{Factoriser au maximum les expressions suivantes
	\begin{enumerate}
		\item $(3x+2)^2-(x-5)^2$
		\item $(x-4)(3x-7)-(x-4)^2-(4-x)(2x+7)$
		\item $(9x^2+12x+4)-2x(3x+2)+(4-9x^2)$
	\end{enumerate}}
	\onslide<2->{}
	
\end{frame}
\end{document}
